\chapter{Extended Listings}
\label{ch:app}

The listings in the main chapters are shortened to the statements that
carry the argument. This appendix reproduces the surrounding functions
at the length that a re-implementer would want.

\section{Periodic diagnostic stack}

\begin{lstlisting}[style=cstyle,caption={Diagnostic register list and Send\_Status sketch.}]
#define CAN_NUM_MSG  10

uint8_t msg_tx_diag_can[CAN_NUM_MSG] = {
    MC_PROTOCOL_REG_UNDEFINED,    /* sequence counter, 100 ms */
    MC_PROTOCOL_REG_SPEED_REF,
    MC_PROTOCOL_REG_TORQUE_REF,
    MC_PROTOCOL_REG_FLUX_REF,
    MC_PROTOCOL_REG_BUS_VOLTAGE,
    MC_PROTOCOL_REG_HEATS_TEMP,
    MC_PROTOCOL_REG_SPEED_MEAS,
    MC_PROTOCOL_REG_TORQUE_MEAS,
    MC_PROTOCOL_REG_FLUX_MEAS,
    MC_PROTOCOL_REG_STATUS
};

uint8_t Send_Status(MCP_Handle_t *pHandle)
{
    if (FCP_TRANSFER_IDLE == pHandle->pFCP->TxFrameState) {
        for (i = 0; i < CAN_NUM_MSG; i++) {
            support_buffer = UI_GetReg(&pHandle->_Super,
                                       msg_tx_diag_can[i], NULL);
            /* pack support_buffer little-endian, ID = reg + 0x40 */
            pHandle->pFCP->TxFrame.Code = code;
            pHandle->pFCP->TxFrame.Size = size;
            DBCCAN_TX_IRQ_Handler(pHandle->pFCP);
        }
    }
    return 0;
}
\end{lstlisting}

\section{Flux-ramp call chain}

\begin{lstlisting}[style=cstyle,caption={Complete UI / MCI / STC flux-ramp chain.}]
__weak bool UI_ExecFluxRamp(UI_Handle_t *pHandle,
                            qd_t Vqd, uint16_t hDurationms)
{
    MCI_Handle_t *pMCI = pHandle->pMCI[pHandle->bSelectedDrive];
    MCI_ExecFluxRamp(pMCI, Vqd, hDurationms);
    return true;
}

__weak void MCI_ExecFluxRamp(MCI_Handle_t *pHandle,
                             qd_t Iqdref, uint16_t hDurationms)
{
    pHandle->lastCommand  = MCI_EXECFLUXTRAMP;
    pHandle->hFinalFlux   = Iqdref.d;
    pHandle->hDurationms  = hDurationms;
    pHandle->CommandState = MCI_COMMAND_NOT_ALREADY_EXECUTED;
}

__weak bool STC_ExecCurrentRamps(SpeednTorqCtrl_Handle_t *pHandle,
                                 int16_t hFinalFlux,
                                 uint16_t hDurationms)
{
    bool     AllowedRange = true;
    uint32_t wAux;
    int32_t  wAux1;
    int16_t  hCurrentReferenceId = STC_GetFluxRef(pHandle);

    if (AllowedRange == true) {
        STC_SetControlMode(pHandle, STC_TORQUE_MODE);
        if (hDurationms == 0u) {
            pHandle->FluxRef             = (int32_t)hFinalFlux * 65536;
            pHandle->RampRemainingStepId = 0u;
            pHandle->IncDecAmountId      = 0;
        } else {
            pHandle->TargetFinalId = hFinalFlux;
            wAux  = (uint32_t)hDurationms
                  * (uint32_t)pHandle->STCFrequencyHz;
            wAux /= 1000u;
            pHandle->RampRemainingStepId = wAux + 1u;
            wAux1  = ((int32_t)hFinalFlux
                    - (int32_t)hCurrentReferenceId) * 65536;
            wAux1 /= (int32_t)pHandle->RampRemainingStepId;
            pHandle->IncDecAmountId = wAux1;
        }
    }
    return AllowedRange;
}
\end{lstlisting}

\section{CAPL command and diagnostic events}

\begin{lstlisting}[style=caplstyle,caption={CAPL command and diagnostic events as deployed.}]
on sysvar SET_SPEED_RAMP::SET_RAMP
{
  message SET_RAMP msgSpeedRamp;
  if (@this == 1 &&
      is_finalSpeed_ready == 1 &&
      is_durSpeed_ready   == 1)
  {
    msgSpeedRamp.FinalSpeed_execrmp = g_finalSpeed_value;
    msgSpeedRamp.Duration_execrmp   = g_durSpeed_value;
    output(msgSpeedRamp);
  }
}

void send_pi_message(word reg_id, int val, int denom)
{
  message SET_REG msg;
  msg.REGISTER_ID          = reg_id;
  msg.REGISTER_VALUE       = val;
  msg.REGISTER_VALUE_DENOM = denom;
  output(msg);
}

on message DIAG_TORQUE_REF
{
  @DIAG_MESSAGES::TORQUE_REF = this.DWord(0);
}

on message REG_VALUE
{
  @GET_REG::K_num = this.GET_K_NUMERATOR;
  @GET_REG::K_den = this.GET_K_DENOMINATOR;
}
\end{lstlisting}
